\exo \textit{[Calculer]}

Entre 2000 et 2017, les coûts du secteur Santé ont augmenté de 40~\% dans l'Union Européenne.

\begin{enumerate}
\item Expliquer pourquoi, si les coûts étaient multipliés par $\dfrac{5}{7}$, ils reviendraient à leur niveau de l'an 2000.
\item Indiquer le taux de cette évolution réciproque, en pourcentage. Arrondir au dixième. 
\end{enumerate}

\exo \textit{[Calculer]}

Le 24 octobre 2018, le NASDAQ, l'un des principaux indices boursiers américains, a baissé de 4,51~\%. Quel aurait dû être son taux d'évolution en pourcentage le 25 octobre pour qu'il retrouve sa valeur du 23 octobre? Arrondir au centième.

\exo \textit{[Calculer]}

On donne ci-dessous un algorithme écrit en langage naturel.

\medskip

\begin{tabular}{|l}
1: Saisir la valeur de $t$\\

2: C $\leftarrow 1/(1+t/100)$ \\

3: t $\leftarrow (C-1)*100$ \\

4: Afficher t\\
\end{tabular}


\begin{enumerate}
\item Que fait cet algorithme?
\item Les lignes 2 et 3 pourraient être remplacées par une seule instruction. Laquelle?
\item Écrire, en langage PYTHON, une fonction taux qui renverrait le même résultat que l'algorithme ci-dessus.
\end{enumerate}

\exo \textit{[Calculer]}

Réduire au même dénominateur.

\medskip

a)~$\dfrac{1}{x}-3$ pour $x \neq 0$ \hspace{1cm} b)~$\dfrac{x+2}{x-3}-1$ pour $x \neq 3$ \hspace{1cm} c)~$\dfrac{7}{3x+1}+2$ pour $x \neq -\dfrac{1}{3}$.


\exo \textit{[Modéliser, Calculer]}

Marion affirme : \og Le nombre $\dfrac{x+5}{x-1}$  est inférieur à 10 lorsque je donne à $x$ des valeurs supérieures à 2. \fg{} 

Qu'en pensez-vous? Justifier.

\medskip

\exo \textit{[Modéliser, Calculer]}

Problème  1 :

Un électricien dispose de deux résistances $R_1$ et $R_2$, l'une fixe de 10~$\Omega$. Il associe ces deux résistances en parallèle. On rappelle que, dans ce cas, la résistance équivalente $R_e$ est telle que $\dfrac{1}{R_e}=\dfrac{1}{R_1}+\dfrac{1}{R_2}$.

On note $x$ la valeur de la résistance variable.

\begin{enumerate}
\item Démontrer que la valeur de la résistance équivalente $R_e$ est donnée par $R_e=\dfrac{10x}{x+10}$.
\item Déterminer les valeurs de $x$ pour lesquelles la valeur de la résistance équivalente est supérieure à 4~$\Omega$.
\end{enumerate}

Problème  2 :

Arthur part de son domicile à vélo pour rejoindre le village voisin où habite son amie. Sa vitesse, à l'aller, est de 40~$km.h^{-1}$. Après sa visite, il revient à son domicile à la vitesse de $x$~$km.h^{-1}$.


\begin{enumerate}
\item Démontrer que la vitesse moyenne d'Arthur sur le trajet aller-retour est donnée par $v(x)=\dfrac{80x}{x+40}$.
\item Déterminer les valeurs de $x$ pour lesquelles la vitesse moyenne est supérieure à 35~$km.h^{-1}$.
\end{enumerate}

Coup de pouce : Introduire la distance d entre le domicile d'Arthur et le village de son amie. Exprimer en fonction de d la durée totale du parcours aller-retour et en déduire la vitesse moyenne sur ce trajet.


\exo \textit{[Chercher, Calculer]}

Dire si l'affirmation est vraie ou fausse, puis justifier.

\og Si Guillaume gagne 25\% de plus que Maud, alors Maud gagne 20\% de moins que Guillaume. \fg{} 




